% !TEX program = xelatex
% Báo cáo phân vùng polyp (tiếng Việt). Biên dịch: latexmk -xelatex Polyp-Segmentation_Report-VI_v1.tex
% Chữ đỏ (\ph, \tbd, \figph) đánh dấu nội dung còn phải viết hoặc đo.
\documentclass[10pt,twoside]{article}
\usepackage{report_style}
\usepackage{polyglossia}
\setdefaultlanguage{vietnamese}
\setotherlanguage{english}
\gappto\captionsvietnamese{%
  \renewcommand{\figurename}{Hình}%
  \renewcommand{\tablename}{Bảng}%
  \renewcommand{\refname}{Tài liệu tham khảo}}

\begin{document}

\makereporttitle
  {Phân vùng polyp đại tràng trên ảnh nội soi: so sánh mô hình CNN và Transformer, cải thiện khả năng tổng quát hóa trên dữ liệu ngoài miền}
  {Nguyễn An Trần Điền\authorand{và}\ph{[Thành viên nhóm]}}
  {Trường Đại học Công nghệ Thông tin,\\ Thành phố Hồ Chí Minh, Việt Nam.}

\begin{reportabstract}[Tóm tắt]
\abslabel{Mục đích} Các mô hình phân vùng polyp đạt Dice gần 0,90 trên ảnh của
những trung tâm có trong dữ liệu huấn luyện, nhưng giảm mạnh trên ảnh từ trung tâm
khác: PraNet đạt 0,898 trên Kvasir-SEG nhưng chỉ 0,628 trên ETIS. Phần lớn các so
sánh đã công bố cũng chỉ dựa vào Dice. Nghiên cứu này so sánh các mô hình CNN và
Transformer trên cùng một pipeline huấn luyện, và kiểm tra xem những thay đổi chi
phí thấp có thu hẹp được khoảng cách ngoài miền hay không.
\abslabel{Phương pháp} Mọi mô hình được huấn luyện trên split PraNet của
Kvasir-SEG và CVC-ClinicDB (1.450 ảnh) và kiểm tra trên năm bộ, trong đó ba bộ đến
từ các trung tâm không xuất hiện khi huấn luyện. Nhóm tái hiện U-Net, MultiResUNet,
CE-Net và TransUNet \ph{[và Swin-Unet]} với cùng dữ liệu, augmentation, hàm mất mát
và lịch huấn luyện, sau đó thêm hàm mất mát chú ý đường biên và augmentation theo
màu sắc. Các mô hình được so sánh bằng Dice, IoU, recall và HD95; chênh lệch được
kiểm định bằng phép kiểm định hạng có dấu Wilcoxon trên điểm của từng ảnh.
\abslabel{Kết quả} \ph{[Placeholder] Dice của từng mô hình trên tập trong miền và
ngoài miền; độ lớn khoảng cách tổng quát hóa; Transformer có thắng CNN một cách có ý
nghĩa thống kê hay không; tác dụng của từng cải tiến trong ablation; tốc độ suy
luận.}
\abslabel{Kết luận} \ph{[Placeholder] Một hoặc hai câu về phát hiện chính, kèm hạn
chế chính (tập huấn luyện nhỏ, chỉ dùng benchmark công khai, ít lần chạy cho mỗi
cấu hình).}
\keywords[Từ khóa]{phân vùng polyp, nội soi đại tràng, U-Net, Vision Transformer,
tổng quát hóa miền, boundary loss, kiểm định thống kê}
\end{reportabstract}

% =====================================================================
\section{Giới thiệu}
Phần lớn ung thư đại trực tràng phát triển từ polyp trong đại tràng, và việc phát
hiện, cắt bỏ polyp khi nội soi giúp giảm nguy cơ ung thư \ph{[trích dẫn]}. Polyp
khác nhau về kích thước, hình dạng và màu sắc, nhiều polyp trông gần giống niêm mạc
xung quanh nên vẫn có thể bị bỏ sót. Phân vùng tự động khoanh vùng polyp đến từng
điểm ảnh và có thể hỗ trợ bác sĩ trong lúc nội soi. Các mạng encoder--decoder dựa trên
U-Net~\cite{ronneberger2015unet} là cách tiếp cận chuẩn, còn các mô hình gần đây mở
rộng hoặc thay thế encoder tích chập bằng
Transformer~\cite{chen2021transunet,cao2022swinunet}.

Điểm số trên các benchmark quen thuộc rất cao, nhưng giảm mạnh trên dữ liệu từ
trung tâm khác. PraNet đạt Dice trung bình 0,898 trên Kvasir-SEG và 0,899 trên
CVC-ClinicDB, hai nguồn của dữ liệu huấn luyện, nhưng chỉ 0,709 trên CVC-ColonDB và
0,628 trên ETIS~\cite{fan2020pranet}. Cách đánh giá còn là vấn đề thứ hai. Một bài
kiểm định 27 công trình phân vùng polyp cho thấy 26 bài không có kiểm định ý nghĩa
thống kê, 25 bài bỏ qua HD95 và 25 bài bỏ qua recall; trên sáu trung tâm chưa thấy
của PolypGen, mọi mô hình được đánh giá lại chỉ còn Dice
0,71--0,76~\cite{urooj2026audit}. Vì vậy, chênh lệch Dice nhỏ đo trên tập trong miền
không cho biết nhiều về việc mô hình nào hoạt động tốt hơn trên dữ liệu mới.

Báo cáo này không đặt mục tiêu đạt kết quả tốt nhất hiện nay. Báo cáo giải quyết một
câu hỏi hẹp hơn: trên cùng một pipeline có kiểm soát, các mô hình CNN và Transformer
trong danh sách bài đọc của môn
học~\cite{ronneberger2015unet,chen2021transunet,cao2022swinunet,ibtehaz2020multiresunet,gu2019cenet,isensee2018nnunet}
so với nhau ra sao trên dữ liệu trong miền và ngoài miền, và những thay đổi chi phí
thấp có thu hẹp được khoảng cách giữa hai loại dữ liệu này không?

Báo cáo trả lời ba câu hỏi nghiên cứu:
\begin{itemize}
  \item \textbf{RQ1.} Với khoảng 1.450 ảnh huấn luyện, các mô hình dựa trên
  Transformer (TransUNet, Swin-Unet) có vượt các mô hình CNN (MultiResUNet, CE-Net)
  một cách có ý nghĩa thống kê không?
  \item \textbf{RQ2.} Mô hình nào giữ hiệu năng tốt nhất trên các tập ngoài miền
  CVC-ColonDB, ETIS và CVC-300?
  \item \textbf{RQ3.} Hàm mất mát chú ý đường biên và augmentation theo màu sắc thu
  hẹp khoảng cách tổng quát hóa được bao nhiêu, và thành phần nào đóng góp chính?
\end{itemize}

Phần còn lại được tổ chức như sau. Mục~\ref{sec:related} điểm lại các mạng CNN
encoder--decoder, phân vùng dựa trên Transformer và các benchmark polyp.
Mục~\ref{sec:problem} phát biểu bài toán và mô tả dữ liệu. Mục~\ref{sec:pipeline}
trình bày pipeline thực nghiệm và các mô hình, tiếp theo là hàm mất mát và các cải
tiến đề xuất ở Mục~\ref{sec:methods}. Mục~\ref{sec:setup} và~\ref{sec:results}
trình bày thiết lập thực nghiệm và kết quả. Mục~\ref{sec:limits} nêu các giới hạn
về độ tin cậy, và Mục~\ref{sec:conclusion} kết luận.

% =====================================================================
\section{Các công trình liên quan}\label{sec:related}
\subsection{Mạng CNN encoder--decoder cho phân vùng ảnh y tế}
U-Net kết hợp một encoder thu hẹp với một decoder mở rộng và chuyển đặc trưng độ
phân giải cao qua các skip connection, nhờ đó học được từ ít ảnh có
nhãn~\cite{ronneberger2015unet}. MultiResUNet thay hai lớp tích chập ở mỗi tầng bằng
một khối MultiRes, gồm chuỗi tích chập $3\times3$ có đầu ra được ghép lại để bao
nhiều vùng tiếp nhận, và thêm các ``Res path'' có kết nối phần dư trên skip connection
để giảm khoảng cách ngữ nghĩa giữa đặc trưng encoder và decoder; mô hình được đánh
giá trên năm bộ dữ liệu y sinh, trong đó có CVC-ClinicDB và
ISIC~2018~\cite{ibtehaz2020multiresunet}. CE-Net đặt một khối tích chập giãn nở dày
đặc (DAC) và một khối gộp đa nhân có phần dư (RMP) giữa encoder ResNet đã huấn luyện
trước và decoder, để thu ngữ cảnh đa tỉ lệ mà không giảm thêm độ phân
giải~\cite{gu2019cenet}. nnU-Net không đề xuất kiến trúc mới: khung này tự cấu hình
tiền xử lý, cấu trúc mạng, huấn luyện và hậu xử lý từ ``dấu vân tay'' của bộ dữ liệu,
và được đánh giá trên mười bài toán của Medical Segmentation
Decathlon~\cite{isensee2018nnunet}.

\subsection{Transformer cho phân vùng ảnh y tế}
Tích chập nắm bắt mẫu cục bộ, và vùng tiếp nhận hạn chế khiến khó mô hình hóa ngữ
cảnh tầm xa. TransUNet đưa bản đồ đặc trưng CNN vào một encoder Vision Transformer
rồi giải mã bằng bộ upsample nhiều tầng cùng skip connection kiểu
U-Net~\cite{chen2021transunet}. Swin-Unet xây toàn bộ mạng hình chữ U từ các khối
Swin Transformer, với cửa sổ dịch chuyển giúp giảm chi phí self-attention, và
upsample bằng các lớp patch expanding~\cite{cao2022swinunet}. Cả hai mô hình được
đánh giá trên bộ CT đa cơ quan Synapse và bộ MRI tim ACDC, không phải trên ảnh nội
soi đại tràng.

\subsection{Benchmark phân vùng polyp và thực hành đánh giá}
PraNet đưa ra giao thức mà phần lớn các công trình về polyp áp dụng: huấn luyện trên
900 ảnh Kvasir-SEG và 550 ảnh CVC-ClinicDB, kiểm tra trên năm
bộ~\cite{fan2020pranet}. PolypGen bổ sung ảnh từ sáu trung tâm để kiểm tra khả năng
tổng quát hóa~\cite{ali2023polypgen}. Bài kiểm định trong~\cite{urooj2026audit} ghi
nhận ít nhất năm giao thức chia dữ liệu không tương thích đang được dùng, lỗi chép số
liệu trong các bảng tổng hợp, và chênh lệch HD95 tới khoảng 16~px giữa các thư viện
tính độ đo. Bài này đề xuất một danh mục báo cáo: HD95 và recall là độ đo chính, nêu
rõ giao thức chia dữ liệu, có tập kiểm tra bên ngoài, và kiểm định ý nghĩa trên điểm
của từng ảnh. Metrics Reloaded đưa ra hướng dẫn chung về chọn độ đo cho bài toán
phân vùng~\cite{maierhein2024metrics}. Báo cáo này làm theo các khuyến nghị đó.

% =====================================================================
\section{Phát biểu bài toán và dữ liệu}\label{sec:problem}
\subsection{Phát biểu bài toán}
Mỗi ảnh $x\in\mathbb{R}^{H\times W\times 3}$ có một mặt nạ nhị phân
$y\in\{0,1\}^{H\times W}$ đánh dấu các điểm ảnh thuộc polyp. Mô hình $f_\theta$ trả
về một bản đồ xác suất, và mặt nạ dự đoán có được bằng cách lấy ngưỡng:
\begin{equation}
  p = f_\theta(x)\in[0,1]^{H\times W},\qquad \hat{y} = \mathbb{1}[\,p \ge 0{,}5\,].
\end{equation}
Ảnh huấn luyện đến từ miền nguồn $\mathcal{D}_S$ (Kvasir-SEG và CVC-ClinicDB). Ảnh
kiểm tra đến từ $\mathcal{D}_S$ và từ $K$ miền đích
$\mathcal{D}_T^1,\dots,\mathcal{D}_T^K$ không bao giờ xuất hiện khi huấn luyện
(CVC-ColonDB, ETIS và CVC-300). Với một điểm số theo ảnh $s$ (Dice trong phân tích
chính), khoảng cách tổng quát hóa của một mô hình là
\begin{equation}\label{eq:gap}
  \Delta = \bar{s}\bigl(\mathcal{D}_S^{\text{test}}\bigr)
  - \frac{1}{K}\sum_{k=1}^{K}\bar{s}\bigl(\mathcal{D}_T^{k}\bigr),
\end{equation}
trong đó $\bar{s}$ là điểm trung bình trên các ảnh của một tập. Mô hình tổng quát hóa
tốt hơn khi đạt điểm cao trên các miền đích và có $\Delta$ nhỏ.

\subsection{Nguồn dữ liệu}
Bảng~\ref{tab:data} liệt kê năm bộ dữ liệu công khai. Kvasir-SEG và CVC-ClinicDB cung
cấp ảnh huấn luyện; CVC-ColonDB, ETIS và CVC-300 chỉ dùng để kiểm tra. Cả năm bộ được
dùng theo dạng đã phát hành kèm PraNet~\cite{fan2020pranet}.

\begin{table}[t]
\centering
\caption{Các bộ dữ liệu và vai trò.}\label{tab:data}
\begin{tabularx}{\linewidth}{@{}lll>{\raggedright\arraybackslash}X@{}}
\toprule
Bộ dữ liệu & Số ảnh & Độ phân giải & Vai trò \\
\midrule
Kvasir-SEG~\cite{jha2020kvasir} & 1.000 & khác nhau & Huấn luyện; kiểm tra trong miền \\
CVC-ClinicDB~\cite{bernal2015clinicdb} & 612 & $384\times288$ & Huấn luyện; kiểm tra trong miền \\
CVC-ColonDB~\cite{bernal2012colondb} & 380 & $574\times500$ & Kiểm tra ngoài miền \\
ETIS-LaribPolypDB~\cite{silva2014etis} & 196 & $1225\times966$ & Kiểm tra ngoài miền \\
CVC-300 (EndoScene)~\cite{vazquez2017endoscene} & 60 & \ph{[kiểm tra]} & Kiểm tra ngoài miền \\
\bottomrule
\end{tabularx}
\end{table}

\subsection{Chia dữ liệu}
Bảng~\ref{tab:split} cho số ảnh của mỗi tập. Ảnh huấn luyện và kiểm tra theo split cố
định của PraNet~\cite{fan2020pranet}. Mười phần trăm ảnh huấn luyện của mỗi bộ nguồn
được tách ngẫu nhiên (seed 42) làm tập validation, để chọn checkpoint mà không đụng
đến tập kiểm tra nào. Vì ba bộ CVC do các nhóm có liên quan thu thập, các ảnh gần
trùng giữa tập huấn luyện và tập kiểm tra được kiểm tra bằng perceptual hashing trước
khi huấn luyện, như trong~\cite{urooj2026audit}. Hình~\ref{fig:samples} minh họa một
số ảnh mẫu.

\begin{table}[t]
\centering
\caption{Số ảnh của mỗi tập.}\label{tab:split}
\begin{tabular}{@{}lrrrrrr@{}}
\toprule
Tập & Kvasir & ClinicDB & ColonDB & ETIS & CVC-300 & Tổng \\
\midrule
Huấn luyện & 810 & 495 & -- & -- & -- & 1.305 \\
Validation &  90 &  55 & -- & -- & -- &   145 \\
Kiểm tra   & 100 &  62 & 380 & 196 & 60 & 798 \\
\bottomrule
\end{tabular}
\end{table}

\begin{figure}[t]
\centering
\figph{[Hình placeholder] Một ảnh và mặt nạ của mỗi bộ trong năm bộ dữ liệu, cho thấy
khác biệt về màu sắc, ánh sáng và độ phân giải. Tạo bằng
scripts/plot\_dataset\_examples.py.}
\caption{Ảnh mẫu và mặt nạ nhãn của từng bộ dữ liệu.}\label{fig:samples}
\end{figure}

% =====================================================================
\section{Pipeline thực nghiệm}\label{sec:pipeline}
Hình~\ref{fig:pipeline} trình bày pipeline, và Bảng~\ref{tab:models} liệt kê các mô
hình. Mọi mô hình đi qua cùng các tầng với cùng dữ liệu, augmentation, hàm mất mát và
lịch huấn luyện; chỉ kiến trúc thay đổi, nên chênh lệch kết quả có thể quy cho kiến
trúc chứ không phải cho cách huấn luyện. Các tầng tô xám là nơi áp dụng các cải tiến
ở Mục~\ref{sec:methods}.

\begin{figure}[t]
\centering
\begin{tikzpicture}[
  every node/.append style={execute at begin node={\hyphenpenalty=10000}},
  box/.style={draw, rounded corners=2pt, align=center, font=\scriptsize,
    text width=72pt, minimum height=40pt, inner sep=2pt},
  hl/.style={box, fill=black!8},
  arr/.style={-{Stealth[length=3.5pt]}, semithick}]
  \node[box] (a) at (40pt,0)    {\textbf{1 Dữ liệu}\\ split PraNet,\\ 5 bộ dữ liệu};
  \node[box] (b) at (136pt,0)   {\textbf{2 Tiền xử lý}\\ đổi kích thước,\\ chuẩn hóa ImageNet};
  \node[hl]  (c) at (232pt,0)   {\textbf{3 Augmentation}\\ lật, xoay, scale,\\ màu sắc (C3)};
  \node[box] (d) at (328pt,0)   {\textbf{4 Mô hình}\\ U-Net, CE-Net,\\ MultiResUNet,\\ TransUNet, Swin-Unet};
  \node[hl]  (e) at (328pt,-56pt) {\textbf{5 Huấn luyện}\\ BCE + Dice,\\ boundary loss (C2)};
  \node[hl]  (f) at (232pt,-56pt) {\textbf{6 Suy luận}\\ ngưỡng 0,5,\\ TTA (C3)};
  \node[box] (g) at (136pt,-56pt) {\textbf{7 Đánh giá}\\ Dice, IoU,\\ recall, HD95};
  \node[box] (h) at (40pt,-56pt)  {\textbf{8 Thống kê, demo}\\ kiểm định Wilcoxon,\\ Friedman; demo Gradio};
  \draw[arr] (a) -- (b); \draw[arr] (b) -- (c); \draw[arr] (c) -- (d);
  \draw[arr] (d) -- (e); \draw[arr] (e) -- (f); \draw[arr] (f) -- (g); \draw[arr] (g) -- (h);
  \node[draw, dashed, rounded corners=2pt, font=\scriptsize, text width=356pt,
    align=center, inner sep=4pt] at (184pt,-102pt)
    {Cùng dữ liệu, lịch huấn luyện và seed cho mọi mô hình \,$\bullet$\, chọn checkpoint
     trên validation \,$\bullet$\, ghi điểm theo từng ảnh cho mọi tập kiểm tra};
\end{tikzpicture}
\caption{Pipeline thực nghiệm. Các tầng tô xám được thay đổi bởi các cải tiến đề xuất
(C2: hàm mất mát chú ý đường biên; C3: augmentation theo miền và augmentation khi
suy luận).}\label{fig:pipeline}
\end{figure}

\begin{table}[t]
\centering
\caption{Các mô hình được so sánh (``\#'' là số thứ tự trong danh sách bài đọc của
môn học).}\label{tab:models}
\begin{tabularx}{\linewidth}{@{}l c >{\raggedright\arraybackslash}X >{\raggedright\arraybackslash}X l@{}}
\toprule
Mô hình & \# & Ý tưởng chính & Vấn đề giải quyết & Ưu tiên \\
\midrule
U-Net~\cite{ronneberger2015unet} & -- & Encoder--decoder với skip connection & Mốc so sánh cơ bản & Bắt buộc \\
MultiResUNet~\cite{ibtehaz2020multiresunet} & 59 & Khối MultiRes; Res path & Polyp nhiều kích thước; khoảng cách ngữ nghĩa encoder--decoder & Bắt buộc \\
CE-Net~\cite{gu2019cenet} & 60 & Khối DAC và RMP ở bottleneck & Mất chi tiết không gian do pooling liên tục & Bắt buộc \\
TransUNet~\cite{chen2021transunet} & 57 & Đặc trưng CNN, encoder ViT, decoder U-Net & Vùng tiếp nhận hạn chế của CNN & Bắt buộc$^{*}$ \\
Swin-Unet~\cite{cao2022swinunet} & 56 & Mạng chữ U gồm các khối Swin Transformer & Ngữ cảnh toàn cục với chi phí thấp hơn ViT & Tùy chọn$^{*}$ \\
nnU-Net~\cite{isensee2018nnunet} & 58 & Pipeline U-Net tự cấu hình & Mốc so sánh mạnh không cần chỉnh tay & Tùy chọn \\
\bottomrule
\multicolumn{5}{@{}l}{$^{*}$Cần làm ít nhất một trong hai mô hình Transformer.}
\end{tabularx}
\end{table}

TransUNet và Swin-Unet mặc định nhận ảnh $224\times224$, còn các mô hình CNN theo
PraNet dùng $352\times352$. Kích thước đầu vào của từng mô hình được ghi ở Phụ
lục~\ref{app:settings}, và dự đoán được đưa về độ phân giải gốc trước khi đánh giá,
nên mọi mô hình được chấm trên cùng một mặt nạ. Mỗi lần chạy ghi cấu hình, seed và
điểm theo từng ảnh vào một nhật ký chung. Mô hình tốt nhất được đưa vào một demo
Gradio~\cite{abid2019gradio}, hiển thị mặt nạ dự đoán chồng lên ảnh người dùng tải
lên.

% =====================================================================
\section{Hàm mất mát và các cải tiến đề xuất}\label{sec:methods}
\subsection{Hàm mất mát phân vùng}
Mọi mô hình được huấn luyện với tổng của binary cross-entropy và soft Dice loss:
\begin{equation}\label{eq:loss}
  \mathcal{L}_{\text{seg}} = \mathcal{L}_{\text{BCE}}(p,y)
  + 1 - \frac{2\sum_i p_i y_i + \epsilon}{\sum_i p_i + \sum_i y_i + \epsilon},
\end{equation}
trong đó $i$ chạy trên các điểm ảnh và $\epsilon$ tránh chia cho 0. Cross-entropy cho
gradient ở mọi điểm ảnh, còn thành phần Dice bù cho sự mất cân bằng giữa vùng polyp
nhỏ và phần nền lớn.

\subsection{Hàm mất mát chú ý đường biên (C2)}
Dice và cross-entropy đếm điểm ảnh, nên một dự đoán phủ đúng polyp nhưng lệch biên vài
điểm ảnh gần như không bị phạt, trong khi HD95 lại nhạy với lỗi này. Theo Kervadec và
cộng sự~\cite{kervadec2021boundary}, nhóm thêm một thành phần có trọng số là bản đồ
khoảng cách có dấu $\phi_G$ tới đường biên của mặt nạ nhãn, âm bên trong polyp và
dương bên ngoài:
\begin{equation}\label{eq:boundary}
  \mathcal{L}_{B} = \frac{1}{|\Omega|}\sum_{q\in\Omega}\phi_G(q)\,p(q),
  \qquad
  \mathcal{L} = \mathcal{L}_{\text{seg}} + \lambda\,\mathcal{L}_{B}.
\end{equation}
Xác suất đặt ở xa bên ngoài vùng thật bị phạt nặng hơn xác suất gần biên. Theo khuyến
nghị trong~\cite{kervadec2021boundary}, $\lambda$ bắt đầu nhỏ và tăng dần trong quá
trình huấn luyện, tới \ph{[giá trị]}. Nếu thành phần này làm huấn luyện mất ổn định,
phương án dự phòng là thêm một nhánh phụ dự đoán bản đồ đường biên.

\subsection{Augmentation theo miền và augmentation khi suy luận (C3)}
Nhóm giả định rằng phần lớn khác biệt giữa các trung tâm đến từ màu sắc, ánh sáng và
độ phân giải của máy nội soi. Vì vậy augmentation cơ bản (lật, xoay, scale, đổi độ
sáng nhẹ) được mở rộng bằng các biến đổi quang học mạnh hơn: độ sáng và độ tương
phản, dịch sắc độ và độ bão hòa, gamma, làm mờ và nén JPEG, cài đặt bằng
Albumentations~\cite{buslaev2020albumentations}. Khi suy luận, dự đoán được lấy trung
bình qua phép lật ngang và lật dọc. Tác dụng được đo bằng thay đổi của khoảng cách
$\Delta$ trong công thức~\eqref{eq:gap}.

\subsection{Encoder huấn luyện trước mạnh hơn (C1, tùy chọn)}
Nếu bảng so sánh baseline hoàn thành trước tuần 4, encoder của mô hình tốt nhất được
thay bằng backbone PVTv2~\cite{wang2022pvtv2} hoặc ConvNeXt~\cite{liu2022convnext} huấn
luyện trước trên ImageNet, còn decoder giữ các Res path của MultiResUNet. Mục tiêu là
kết hợp ngữ cảnh toàn cục với chi tiết đa tỉ lệ.

% =====================================================================
\section{Thiết lập thực nghiệm}\label{sec:setup}
\begin{table}[t]
\centering
\caption{Cấu hình thực nghiệm chính.}\label{tab:config}
\begin{tabularx}{\linewidth}{@{}l >{\raggedright\arraybackslash}X@{}}
\toprule
Thành phần & Cấu hình \\
\midrule
Dữ liệu huấn luyện & 1.305 ảnh (tập huấn luyện của split PraNet trừ 10\% làm validation) \\
Kích thước đầu vào & $352\times352$ cho CNN; $224\times224$ cho TransUNet và Swin-Unet \ph{[xác nhận]} \\
Chuẩn hóa & Trung bình và độ lệch chuẩn của ImageNet \\
Augmentation & Lật, xoay, scale ngẫu nhiên, đổi màu nhẹ; bộ biến đổi quang học mạnh hơn cho C3 \\
Bộ tối ưu & AdamW~\cite{loshchilov2019adamw}, learning rate $10^{-4}$, weight decay \ph{[giá trị]} \\
Batch, số epoch & \ph{8--16}, 100 epoch, mixed precision \\
Hàm mất mát & BCE + Dice (công thức~\ref{eq:loss}); thêm boundary loss cho C2 (công thức~\ref{eq:boundary}) \\
Chọn checkpoint & Dice cao nhất trên validation; không dùng tập kiểm tra để chọn \\
Framework, phần cứng & PyTorch~\cite{paszke2019pytorch} \ph{[phiên bản]}; GPU Google Colab hoặc Kaggle \ph{[loại]} \\
Seed & 42 \ph{[thêm hai seed nếu đủ thời gian GPU]} \\
\bottomrule
\end{tabularx}
\end{table}

Bảng~\ref{tab:config} tóm tắt cấu hình. Với mặt nạ dự đoán $\hat{y}$ và mặt nạ nhãn
$y$, các độ đo chồng lấp và độ nhạy là
\begin{equation}
  \text{Dice} = \frac{2|\hat{y}\cap y|}{|\hat{y}|+|y|},\qquad
  \text{IoU} = \frac{|\hat{y}\cap y|}{|\hat{y}\cup y|},\qquad
  \text{Recall} = \frac{|\hat{y}\cap y|}{|y|}.
\end{equation}
Recall được báo cáo vì bỏ sót polyp nguy hiểm hơn báo nhầm. Độ chính xác đường biên
được đo bằng phân vị thứ 95 của khoảng cách Hausdorff đối xứng,
\begin{equation}
  \text{HD95}(\hat{y},y) = \max\Bigl\{ P_{95}\bigl\{d(a,\partial y) : a\in\partial\hat{y}\bigr\},\;
  P_{95}\bigl\{d(b,\partial\hat{y}) : b\in\partial y\bigr\}\Bigr\},
\end{equation}
trong đó $\partial$ là đường biên của mặt nạ và $d$ là khoảng cách Euclid tính bằng
điểm ảnh ở độ phân giải gốc. Vì giá trị HD95 khác nhau giữa các thư
viện~\cite{urooj2026audit}, nhóm dùng một thư viện duy nhất (\ph{[tên, phiên bản]})
cho mọi mô hình, và gán HD95 bằng đường chéo ảnh khi dự đoán rỗng. Mọi độ đo được tính
theo từng ảnh rồi lấy trung bình trên mỗi tập kiểm tra (mDice, mIoU).

Hai mô hình được so sánh bằng kiểm định hạng có dấu Wilcoxon~\cite{wilcoxon1945} trên
cặp điểm Dice của từng ảnh; chênh lệch chỉ được coi là có ý nghĩa khi $p<0{,}05$. Để so
sánh mọi mô hình cùng lúc, nhóm dùng kiểm định Friedman~\cite{friedman1937} rồi kiểm
định hậu nghiệm, theo khuyến nghị của Dem\v{s}ar~\cite{demsar2006}. Giá trị trung bình
trên các tập ngoài miền được báo cáo kèm khoảng tin cậy bootstrap 95\%. Hiệu quả tính
toán được báo cáo bằng số tham số, FLOPs và số khung hình mỗi giây với batch bằng 1.

% =====================================================================
\section{Kết quả và thảo luận}\label{sec:results}
\subsection{So sánh CNN và Transformer (RQ1)}
\ph{[Placeholder] So sánh mDice của mọi mô hình trên năm tập kiểm tra
(Bảng~\ref{tab:main}). Nêu rõ Transformer tốt nhất có tốt hơn CNN tốt nhất một cách có
ý nghĩa thống kê không, và thảo luận vai trò của trọng số huấn luyện trước trên
ImageNet.}

\begin{table}[t]
\centering
\caption{mDice trên năm tập kiểm tra. Số liệu của PraNet lấy từ~\cite{fan2020pranet},
chỉ dùng để đối chiếu.}\label{tab:main}
\begin{tabular}{@{}lcccccc@{}}
\toprule
 & \multicolumn{2}{c}{Trong miền} & \multicolumn{3}{c}{Ngoài miền} & \\
\cmidrule(lr){2-3}\cmidrule(lr){4-6}
Mô hình & Kvasir & ClinicDB & ColonDB & ETIS & CVC-300 & $p$ so với CNN \\
\midrule
PraNet~\cite{fan2020pranet} & 0,898 & 0,899 & 0,709 & 0,628 & 0,871 & -- \\
\midrule
U-Net        & \tbd & \tbd & \tbd & \tbd & \tbd & \tbd \\
MultiResUNet & \tbd & \tbd & \tbd & \tbd & \tbd & \tbd \\
CE-Net       & \tbd & \tbd & \tbd & \tbd & \tbd & \tbd \\
TransUNet    & \tbd & \tbd & \tbd & \tbd & \tbd & \tbd \\
Swin-Unet    & \tbd & \tbd & \tbd & \tbd & \tbd & \tbd \\
\bottomrule
\end{tabular}
\end{table}

\subsection{Tổng quát hóa ngoài miền (RQ2)}
\ph{[Placeholder] Báo cáo khoảng cách $\Delta$, recall và HD95 trên các tập ngoài miền
(Bảng~\ref{tab:ood}). Chỉ ra mô hình giảm ít nhất và thứ hạng có thay đổi giữa Dice và
HD95 hay không (Hình~\ref{fig:gap}).}

\begin{table}[t]
\centering
\caption{Khả năng tổng quát hóa của từng mô hình. Trong miền: Kvasir và ClinicDB; ngoài
miền: ColonDB, ETIS và CVC-300 (trung bình của ba tập, khoảng tin cậy 95\% trong
ngoặc).}\label{tab:ood}
\setlength{\tabcolsep}{5pt}
\begin{tabular}{@{}lccccc@{}}
\toprule
Mô hình & \makecell{mDice\\trong miền} & \makecell{mDice\\ngoài miền} & $\Delta$ & \makecell{Recall\\ngoài miền} & \makecell{HD95 ngoài\\miền (px)} \\
\midrule
U-Net        & \tbd & \tbd & \tbd & \tbd & \tbd \\
MultiResUNet & \tbd & \tbd & \tbd & \tbd & \tbd \\
CE-Net       & \tbd & \tbd & \tbd & \tbd & \tbd \\
TransUNet    & \tbd & \tbd & \tbd & \tbd & \tbd \\
Swin-Unet    & \tbd & \tbd & \tbd & \tbd & \tbd \\
\bottomrule
\end{tabular}
\end{table}

\begin{figure}[t]
\centering
\figph{[Hình placeholder] Trái: mDice trong miền và ngoài miền của từng mô hình (cặp
điểm nối nhau). Phải: phân bố Dice theo từng ảnh trên ETIS.}
\caption{Khoảng cách tổng quát hóa của từng mô hình.}\label{fig:gap}
\end{figure}

\subsection{Ablation các cải tiến đề xuất (RQ3)}
\ph{[Placeholder] Xuất phát từ baseline tốt nhất sau tuần 4, thêm từng thành phần và báo
cáo tác dụng (Bảng~\ref{tab:ablation}). Cải tiến không có ý nghĩa thống kê vẫn được báo
cáo như một kết quả âm.}

\begin{table}[t]
\centering
\caption{Ablation trên mô hình baseline tốt nhất. Giá trị $p$ so sánh từng dòng với
baseline trên Dice ngoài miền theo từng ảnh.}\label{tab:ablation}
\setlength{\tabcolsep}{5pt}
\begin{tabular}{@{}lccccc@{}}
\toprule
Cấu hình & \makecell{mDice\\trong miền} & \makecell{mDice\\ngoài miền} & $\Delta$ & \makecell{HD95 ngoài\\miền (px)} & $p$ \\
\midrule
Baseline tốt nhất \ph{[tên]}       & \tbd & \tbd & \tbd & \tbd & -- \\
+ C3 (augmentation miền + TTA)     & \tbd & \tbd & \tbd & \tbd & \tbd \\
+ C2 (boundary loss)               & \tbd & \tbd & \tbd & \tbd & \tbd \\
+ C2 + C3                          & \tbd & \tbd & \tbd & \tbd & \tbd \\
C1 + C2 + C3 (tùy chọn)            & \tbd & \tbd & \tbd & \tbd & \tbd \\
\bottomrule
\end{tabular}
\end{table}

\subsection{Phân tích lỗi và hiệu quả tính toán}
\ph{[Placeholder] Hiển thị các ảnh có Dice thấp nhất của mỗi tập kiểm tra
(Hình~\ref{fig:errors}), chia kết quả theo kích thước polyp, và so sánh chi phí mô hình
(Bảng~\ref{tab:cost}).}

\begin{table}[t]
\centering
\caption{Kích thước mô hình và tốc độ suy luận (batch bằng 1, \ph{[GPU]}).}\label{tab:cost}
\begin{tabular}{@{}lccc@{}}
\toprule
Mô hình & Số tham số (triệu) & FLOPs (G) & Khung hình/giây \\
\midrule
U-Net        & \tbd & \tbd & \tbd \\
MultiResUNet & \tbd & \tbd & \tbd \\
CE-Net       & \tbd & \tbd & \tbd \\
TransUNet    & \tbd & \tbd & \tbd \\
Swin-Unet    & \tbd & \tbd & \tbd \\
Mô hình đề xuất & \tbd & \tbd & \tbd \\
\bottomrule
\end{tabular}
\end{table}

\begin{figure}[t]
\centering
\figph{[Hình placeholder] Ảnh, mặt nạ nhãn và dự đoán của từng mô hình cho ba ảnh có Dice
thấp nhất của ETIS và CVC-ColonDB.}
\caption{Các trường hợp lỗi điển hình trên ảnh ngoài miền.}\label{fig:errors}
\end{figure}

% =====================================================================
\section{Độ tin cậy và hạn chế}\label{sec:limits}
Bằng chứng của báo cáo bị giới hạn ở một số điểm. Các mô hình được huấn luyện trên
1.305 ảnh từ hai nguồn, còn các tập ngoài miền nhỏ (60--380 ảnh), nên giá trị trung bình
có khoảng tin cậy rộng; chênh lệch nhỏ hơn khoảng này không được diễn giải. Mọi ảnh
kiểm tra đều có ít nhất một polyp, nên báo động nhầm trên khung hình bình thường không
được đo, và ảnh tĩnh do nhóm tạo bộ dữ liệu chọn lọc sạch hơn video thực tế, vốn có
nhòe, phản chiếu ánh sáng và chuyển động.

Thời gian GPU miễn phí giới hạn số lần chạy, nên một số cấu hình chỉ được huấn luyện với
một seed. Siêu tham số được dùng chung chứ không tinh chỉnh riêng cho từng mô hình, điều
này có thể có lợi cho mô hình có công thức huấn luyện gốc gần với công thức chung. Các mô
hình tái hiện được huấn luyện trên dữ liệu và kích thước đầu vào khác với bài báo gốc, nên
điểm số chỉ được so sánh trong pipeline này và với PraNet, không so trực tiếp với số liệu
của bài gốc. Đây là một so sánh học thuật trên dữ liệu công khai; nghiên cứu không có
kiểm chứng lâm sàng, và demo không phải là thiết bị y tế.

% =====================================================================
\section{Kết luận và hướng phát triển}\label{sec:conclusion}
\ph{[Placeholder] Tóm tắt câu trả lời cho RQ1--RQ3 trong một đoạn, kèm các số liệu
chính.}

Hướng phát triển gần nhất là kiểm tra các mô hình trên PolypGen~\cite{ali2023polypgen},
với sáu trung tâm cho đánh giá bên ngoài chặt chẽ hơn các tập benchmark công khai. Đánh
giá trên video nội soi sẽ đo được độ ổn định giữa các khung hình và báo động nhầm trên
khung hình không có polyp. Mô hình nhẹ hơn cho suy luận thời gian thực, và huấn luyện bán
giám sát với khung hình nội soi chưa gán nhãn, là các hướng mở rộng tiếp theo.

\runin{Dữ liệu và mã nguồn.} Mã nguồn, cấu hình, danh sách chia dữ liệu, điểm theo từng
ảnh và script vẽ hình được lưu trong kho mã của đề tài \ph{[URL]}. Các bộ dữ liệu là dữ
liệu công khai, tuân theo điều khoản gốc của từng bộ, và không được phát hành lại.

% =====================================================================
\startappendix[Phụ lục]
\section{Thiết lập riêng của từng mô hình}\label{app:settings}
\begin{table}[h]
\centering
\caption{Các thiết lập khác nhau giữa các mô hình (IN: ImageNet). Các thiết lập còn lại
theo Bảng~\ref{tab:config}.}\label{tab:settings}
\setlength{\tabcolsep}{4pt}
\begin{tabular}{@{}lccccc@{}}
\toprule
Mô hình & \makecell{Kích thước\\đầu vào} & \makecell{Trọng số\\huấn luyện trước} & Batch & Epoch & \makecell{Số tham số\\(triệu)} \\
\midrule
U-Net        & $352\times352$ & không & \tbd & \tbd & \tbd \\
MultiResUNet & $352\times352$ & không & \tbd & \tbd & \tbd \\
CE-Net       & $352\times352$ & ResNet-34 (IN-1k) & \tbd & \tbd & \tbd \\
TransUNet    & $224\times224$ & R50-ViT-B/16 (IN-21k) & \tbd & \tbd & \tbd \\
Swin-Unet    & $224\times224$ & Swin-T (IN-1k) & \tbd & \tbd & \tbd \\
Mô hình đề xuất & \tbd & \tbd & \tbd & \tbd & \tbd \\
\bottomrule
\end{tabular}
\end{table}

\clearpage
\section{Bản đồ tài nguyên tái lập}\label{app:artifacts}
\begin{table}[h]
\centering
\caption{Các tài nguyên chính trong kho mã của đề tài (bố cục dự kiến; cập nhật theo
đường dẫn cuối cùng).}\label{tab:artifacts}
\begin{tabular}{@{}ll@{}}
\toprule
Mục đích & Tài nguyên \\
\midrule
Hướng dẫn chạy           & README.md \\
Danh sách chia dữ liệu   & data/splits/\{train,val,test\}\_*.txt \\
Kiểm tra ảnh trùng       & scripts/check\_duplicates.py \\
Cấu hình thí nghiệm      & configs/*.yaml \\
Mã các mô hình           & src/models/ \\
Hàm mất mát              & src/losses.py \\
Huấn luyện và đánh giá   & scripts/train.py, scripts/evaluate.py \\
Điểm theo từng ảnh       & outputs/metrics/per\_image\_*.csv \\
Kiểm định thống kê       & scripts/stats\_tests.py \\
Hình trong báo cáo       & \makecell[l]{scripts/plot\_report\_figures.py\\ $\rightarrow$ report/figures/} \\
Demo                     & app/gradio\_app.py \\
\bottomrule
\end{tabular}
\end{table}

\printreferences

\end{document}
